Let $F$ be a field. By a curve defined over $F$ we mean a geometrically
irreducible, projective variety of dimension $1$ defined over $F$. Let $C$ be
a smooth curve of genus at least $2$ defined over a number field $F$. As was
conjectured by Mordell and proved by Faltings~\cite{Faltings:ES},
$C(F)$, the set of $F$-rational points of $C$, is finite.

We let $\jac(C)$ denote the Jacobian of $C$. Recall
that $\jac(C)(F)$ is a finitely generated abelian group by the
Mordell--Weil Theorem. 

The aim of this paper is to bound $\# C(F)$ from above. Here is our
first result. 

\begin{theorem}\label{ThmBdRatIntro}
  Let $g \ge 2$ and $d \ge 1$ be integers. Then there exists a constant
  $c = c(g, d) \ge 1$ with the following property. If $C$ is a smooth curve
  of genus $g$  defined
  over a number field $F$  with $[F:\IQ] \le d$, then
  \begin{equation}
    \label{eq:CKexpinrank}
      \#C(F) \le c^{1+\rho}
  \end{equation}
  where $\rho$ is the rank of  $\mathrm{Jac}(C)(F)$.% Moreover $c(g,d)$ depends polynomially on $d$.
\end{theorem}

This theorem gives an affirmative answer to a question posed by
Mazur~\cite[Page~223]{Mazur:00}.
% when the objects in question are defined over the field of algebraic numbers.
See also~\cite[top of
page 234]{mazur1986arithmetic} for an earlier question. Before this,
 Lang formulated a related  conjecture~\cite[page
 140]{Lang:EllipticCurves78} on the number of integral
points of elliptic curves.

The method of our theorem builds up on the work of many others. At the
core we follow Vojta's proof~\cite{Vojta:siegelcompact} of the
Mordell Conjecture.
%which also provided a bound for $\#C(F)$. 
Vojta's proof was later simplified  by
Bombieri~\cite{Bombieri:Mordell} and further developed by
Faltings~\cite{Faltings:DAAV}. Silverman~\cite{Silverman:twists}
proved a bound of the quality (\ref{eq:CKexpinrank}) if $C$ ranges
over twists of a given smooth curve. 
The bound by de
Diego~\cite{deDiego:97} %for a family of smooth curves 
is of the form $c(g)
7^{\rho}$, where $c(g)>0$ depends only on $g$; the value $7$ had already
arisen in Bombieri's work. But she only counts points
whose height is large in terms of a height of $C$.
Work of David--Philippon \cite{DPvarabII} and R\'emond~\cite{Remond:Decompte}
led to explicit estimates. 
Recently, Alpoge \cite{AlpogeRatPt} \cite[Theorem~6.1.1]{AlpogeThesis} improved  $7$
to $1.872$ and, for $g$ large enough, even to $1.311$. 

% Our approach to Theorem~\ref{ThmBdRatIntro} follows the original
% blueprint towards Mordell's Conjecture laid out by Vojta. [PH]:
% removed because it repeats the above.
On combining the Vojta and Mumford Inequalities
one gets an upper bound for the number of \textit{large points} in
$C(F)$; these are points whose height is sufficiently large relative
to a suitable height of $C$. A lower bound for the N\'eron--Tate
height, such as proved by David--Philippon \cite{DPvarabII}, can be used to
count the number of remaining points which we sometimes call
\textit{small points}. 
Indeed, R\'emond~\cite{Remond:Decompte} made the Vojta and Mumford
Inequalities explicit and obtained 
explicit upper bounds for the number of rational points on curves
embedded in abelian varieties. 
The resulting cardinality bounds
depend on a suitable notion of height of $C$, an artifact of the lower
bounds for the N\'eron--Tate height. 
Later, David--Philippon~\cite{DaPh:07} proved stronger height lower
bounds
in a power of an
elliptic curve. They then 
 obtained uniform estimates
of the quality (\ref{eq:CKexpinrank}) for a curve in a power of
elliptic curves,  thus providing evidence that Mazur's Question had a
positive answer, see also
David--Nakamaye--Philippon's work~\cite{DaNaPh:07}.

We give an overview of the general method  in more detail in
$\mathsection$\ref{subsec:ntdist} below. 
% The main
% innovation of this paper is to get equally uniform bounds for the
% number of points on $C$ whose height is not sufficiently large.

The main innovation of this paper is to prove a lower bound for the
N\'eron--Tate height that is sufficiently strong to eliminate the
dependency on the height of $C$. This leads to a uniform estimate as
in Theorem \ref{ThmBdRatIntro}. In prior work~\cite{DGH1p} we
applied the earlier height lower bound~\cite{GaoHab}
to recover a variant of Theorem~\ref{ThmBdRatIntro} in
a one-parameter family of smooth curves.

We now explain some further  results that follow from the approach
described above. 
For an integer $g\ge 1$, let $\mathbb{A}_{g,1}$ denote the coarse moduli
space of principally polarized abelian varieties of dimension $g$.
This is an irreducible quasi-projective variety which we can take as
defined over $\IQbar$, the algebraic closure of $\IQ$ in $\IC$.
Suppose we are presented with an immersion
$\iota\colon \mathbb{A}_{g,1}\rightarrow \IP^m_\IQbar$ into projective
space.
Let $h\colon \IP^m_\IQbar(\IQbar)\rightarrow\IR$ denote the absolute
logarithmic Weil height,
\textit{cf.}~\cite[$\mathsection$1.5.1]{BG}.
For brevity, we sometimes call $h$ the Weil height.
If $C$ is a smooth curve of genus $g\ge 2$ defined over $\IQbar$ and if
$P_0\in C(\IQbar)$, then we can consider $C-P_0$ as a curve in
$\mathrm{Jac}(C)$ via the Abel--Jacobi map.
We use $[\mathrm{Jac}(C)]$ to denote the point in $ \mathbb{A}_{g,1}(\IQbar)$  parametrizing $\mathrm{Jac}(C)$.

% Let $A$ be an abelian variety defined over $\IQbar$
An abelian group $\Gamma$ is said to have finite rank if
$\Gamma\otimes\IQ$ is a finite dimensional $\IQ$-vector space. In this
case $\dim\Gamma\otimes\IQ$ is the rank of $\Gamma$. 
Consider an abelian variety $A$ defined over $\IC$ and let $\Gamma$ be
a finite rank subgroup of $A(\IC)$. Lang~\cite{Lang:Division}
conjectured that a
curve in $A$ intersects $\Gamma$ in a finite set
unless the curve is smooth of genus $1$. The Conjecture follows from
Faltings's Theorem~\cite{Faltings:ES} and work of
Raynaud~\cite{Raynaud:MMcurves}.

%What follows is a more precise version of Theorem~\ref{ThmBdRatIntro}.
% Theorem~\ref{ThmBdRatIntro} follows from the following theorem,
The following theorem is more in the spirit of \cite[top of page
234]{mazur1986arithmetic}.

\begin{theorem}
\label{ThmBdFinRank}
%% \Gcomment{Separated the two constants. This may help if later on we remove the ``large height'' hypothesis.}
%% PH: Good, I renamed them to fit our convention on constants
Let $g \ge 2$ and let $\iota$ be as above. Then there exist two
constants $c_1= c_1(g,\iota) \ge 0$ and $c_2 = c_2(g,\iota) \ge 1$ with the following property. Let $C$
be a  smooth curve of genus $g$  defined over $\IQbar$,
let $P_0 \in C(\overline \IQ)$, and let $\Gamma$ be a
subgroup of $\mathrm{Jac}(C)(\overline \IQ)$ of finite rank $\rho\ge
0$. If $h(\iota([\mathrm{Jac}(C)])) \ge c_1$, then
\[
\#(C(\overline \IQ)-P_0) \cap \Gamma  \le c_2^{1+\rho}.
\]
\end{theorem}

The following corollary follows from Theorem \ref{ThmBdFinRank} applied to 
$\Gamma = \mathrm{Jac}(C)(\IQbar)_{\mathrm{tors}}$, which has rank $0$. 


\begin{corollary}
%\label{CorBdAlgTorIntro}
Let $g \ge 2$ and let $\iota$ be as above.
Then there exist two
constants $c_1= c_1(g,\iota) \ge 0$ and $c_2 = c_2(g,\iota) \ge 1$ with the following property. Let $C$
be a  smooth curve of genus $g$  defined over $\overline \IQ$ and  let $P_0 \in C(\overline \IQ)$.
If $h(\iota([\mathrm{Jac}(C)])) \ge c_1$, then
\[
\# (C(\overline \IQ)-P_0) \cap \mathrm{Jac}(C)(\IQbar)_{\mathrm{tors}}  \le c_2.
\]
\end{corollary}

As in Theorem \ref{ThmBdRatIntro} we can drop the condition on the
height of the Jacobian by working over a number field of bounded
degree. 

\begin{theorem}\label{thm:BdTorIntroNF}
  Let $g \ge 2$ and $d \ge 1$ be integers. Then there exists a constant
  $c = c(g, d) \ge 1$ with the following property. Let $C$ be a 
  smooth curve of genus $g$ defined over a
  number field $F\subset\IQbar$ 
  with $[F:\IQ] \le d$ and let $P_0 \in C(\IQbar)$, then
  \begin{equation*}
    \# (C(\overline \IQ)-P_0) \cap \mathrm{Jac}(C)(\IQbar)_{\mathrm{tors}} \le c.
  \end{equation*}
\end{theorem}

Let us recall some previous results towards Mazur's Question for
rational points, \textit{i.e.}, towards Theorem~\ref{ThmBdRatIntro}.
Based on the method of Vojta, Alpoge~\cite{AlpogeRatPt} proved that
the average number of rational points on a curve of genus $2$ with a
marked Weierstrass point is bounded. Let $C$ be a smooth curve of
genus $g\ge 2$ defined
over a number field $F\subset\IQbar$. The Chabauty--Coleman
approach~\cite{Chabauty, Coleman:effCha} yields estimates under an additional
hypothesis on the rank of Mordell--Weil group. For example, if
$\jac(C)(F)$ has rank at most $g-3$, Stoll~\cite{Stoll:Uniform} showed
that $\#C(F)$ is bounded solely in terms of $[F:\IQ]$ and $g$ if $C$ is
hyperelliptic; Katz--Rabinoff--Zureick-Brown~\cite{KatzRabinoffZB}
later, under the same rank hypothesis, removed the hyperelliptic
hypothesis. Checcoli,
Veneziano, and Viada~\cite{CVV:17} obtain an effective height bound
under a restriction on 
 the Mordell--Weil rank.

As for algebraic torsion points, \textit{i.e.}, in the direction of  Theorem~\ref{thm:BdTorIntroNF}, DeMarco--Krieger--Ye \cite{DeMarcoKriegerYeUniManinMumford} proved a bound on the cardinality of torsion points on any genus $2$
curve that admits a degree-two map to an elliptic curve when the
Abel--Jacobi map is based at a Weierstrass point. Moreover, their bound is
independent of $[F:\IQ]$.


\subsection{N\'{e}ron--Tate distance of algebraic points on curves}
\label{subsec:ntdist}
Let $C$ be a smooth curve defined over $\overline \IQ$ of genus $g \ge
2$, let $P_0 \in C(\overline \IQ)$, and let $\Gamma$ be a subgroup of
$\jac(C)(\overline \IQ)$ of finite rank $\rho$. For simplicity we identify $C$ with its image under the Abel--Jacobi embedding $C \rightarrow \jac(C)$ via $P_0$. 

Our proof of Theorem~\ref{ThmBdFinRank} follows the spirit of the
method of Vojta, later generalized by  Faltings.
Let $\hat h\colon \jac(C)(\IQbar)\rightarrow \IR$ denote the
N\'eron--Tate height attached to a symmetric and ample line bundle on $\jac(C)$. 
We divide $C(\overline \IQ) \cap \Gamma$ into two parts:
\begin{itemize}
\item Small points $\left\{P \in C(\overline \IQ) \cap \Gamma
: \hat{h}(P) \le B(C)  \right\}$;
\item Large points $\left\{P \in C(\overline \IQ) \cap \Gamma : \hat{h}(P) > B(C) \right\}$
\end{itemize}
where $B(C)$ is allowed to depend on a suitable height of $C$. It
turns out that
we can take $B(C)=c_0 \max\{1,h(\iota([\jac(C)]))\}$
for some $c_0 = c_0(g,\iota) > 0$. The constant $c_0$ is chosen in a
way that accommodates  both the \textit{Mumford inequality} and
the \textit{Vojta inequality}. Combining these
two inequalities yields an upper bound on the number of large points
by $c_1(g)^{1+\rho}$, see for example Vojta's~\cite[Theorem
6.1]{Vojta:siegelcompact} in the important case where $\Gamma$ is the
group of points of $\jac(C)$ rational over
 a number field or more generally in the work of
 David--Philippon~\cite{DPvarabII,DaPh:07} and R\'emond~\cite{Remond:Decompte}.



Thus in order to prove Theorem~\ref{ThmBdFinRank}, it suffices to bound the number of small points.
In this paper we find such a  bound by studying the
N\'{e}ron--Tate distance of points in $C(\overline \IQ)$.

Roughly speaking,
we find positive constants $c_1,c_2,c_3,$ and $c_4$  that depend on
$g$ and $\iota$, but not 
on $C$, such that  if $h(\iota([\jac(C)]))\ge
c_1$ then for all $P\in C(\IQbar)$ we have the following  alternative.
\begin{itemize}
\label{eq:alternative}
\item Either $P$ lies in  a subset of $C(\overline \IQ)$ of cardinality at most $c_2$,
\item or $\left\{Q \in C(\overline \IQ) : \hat{h}(Q-P) \le h(\iota([\jac(C)]))/c_3 \right\} < c_4$.
\end{itemize}


This dichotomy is stated in Proposition~\ref{PropAlgPtFar}.
In this paper, we make the statement precise by referring to the universal family of genus $g$ smooth curves with suitable level structure, and the N\'eron--Tate height on each Jacobian attached to the tautological line bundle. 
The setup is done in $\mathsection$\ref{SectionSettingUp}.


This proposition can be seen as a relative version of the Bogomolov
conjecture for abelian varieties with large height. It has the
following upshot: If $h(\iota([\jac(C)])) \ge c_1$, then the small
points in $C(\overline \IQ) \cap \Gamma$ lie in a set of uniformly
bounded cardinality, or are contained in $(1+c_0c_3)^\rho$ balls in the
N\'eron--Tate metric, with each ball containing at most $c_4$ points.
This will yield the desired bound in Theorem~\ref{ThmBdFinRank}, as
executed in $\mathsection$\ref{SectionRatPt}.

\subsection{Height inequality and non-degeneracy}\label{subsec:hgtineqnondeg}
We follow the framework presented in our previous work \cite{DGH1p}.
In \textit{loc.cit.} we proved the result for $1$-parameter families,
as an application of the second- and third-named authors' height
inequality \cite[Theorem~1.4]{GaoHab}. Passing to general cases
requires generalizing this height inequality to higher dimensional
bases. The generalization has two parts:  generalizing the inequality
itself under the non-degeneracy condition and generalizing the
criterion of non-degenerate subvarieties.
We execute the first part in the current paper while the second part was 
done by the second-named author in \cite{GaoBettiRank}. Let us explain the setup.

Let $k$ be an algebraically closed subfield of $\IC$. Let $S$ be a
regular, irreducible, quasi-projective variety defined over $k$ that
is Zariski open in an irreducible projective variety
$\overline S \subset\IP_k^m$. Let $\pi \colon \cA \rightarrow S$ be an
abelian scheme of relative dimension $g\ge 1$. We suppose that we are
presented with a closed immersion $\cA\rightarrow\IP_k^n\times S$ over
$S$. On the generic fiber of $\pi$ we assume that this immersion comes
from a basis of the global sections of the $l$-th power of a symmetric
and ample line bundle with $l\ge 4$. If $k=\IQbar$ and as described in
$\mathsection\ref{SubsectionHeight}$ we obtain two height functions,
the restriction of the Weil height
$h\colon \overline{S}(\IQbar)\rightarrow \IR$ and the N\'eron--Tate
height $\hat h_{\cA} \colon \cA(\IQbar)\rightarrow \IR$.

Let  $\ell\ge \bestlevel$ be an integer. Throughout the whole paper, by \textit{level-$\ell$-structure} we always mean \textit{symplectic level-$\ell$-structure}. For the purpose of our main applications, including
Theorems~\ref{ThmBdRatIntro} and \ref{ThmBdFinRank}, it suffices to
work under the following hypothesis. 

\begin{center}
  \label{def:hyp}
  {\tt (Hyp):} $\cA\rightarrow S$ carries a principal polarization and has level-$\ell$-structure for some $\ell \ge 3$.
\end{center}
So in the main body of the paper, we will focus on the case {\tt
(Hyp)}. The general case where {\tt (Hyp)} is not assumed will be
handled in Appendix~\ref{AppHtIneqApp}.


The non-degenerate subvarieties of $\cA$ are defined using the
\textit{Betti map} which we briefly describe here; 
the precise definition will be given by
Proposition~\ref{PropBettiMapApp} and in Proposition~\ref{PropBettiMap} under {\tt (Hyp)}.

For any $s \in S(\IC)$, there exists an open neighborhood $\Delta
\subseteq S^{\mathrm{an}}$ of $s$ which we may assume is simply-connected. Then one can define a basis $\omega_1(s),\ldots,\omega_{2g}(s)$ of the period lattice of each fiber $s \in \Delta$ as holomorphic functions of $s$. Now each fiber $\cA_s = \pi^{-1}(s)$ can be identified with the complex torus $\IC^g/(\IZ \omega_1(s)\oplus \cdots \oplus \IZ\omega_{2g}(s))$, and each point $x \in \cA_s(\IC)$ can be expressed as the class of $\sum_{i=1}^{2g}b_i(x) \omega_i(s)$ for real numbers $b_1(x),\ldots,b_{2g}(x)$. Then $b_{\Delta}(x)$ is defined to be the class of the $2g$-tuple $(b_1(x),\ldots,b_{2g}(x)) \in \IR^{2g}$ modulo $\IZ^{2g}$.
We obtain with a real-analytic map
\[
b_{\Delta} \colon \cA_{\Delta} = \pi^{-1}(\Delta) \rightarrow \mathbb{T}^{2g},
\]
which is fiberwise a group isomorphism and
where $\mathbb{T}^{2g}$ is the real torus of dimension $2g$.


\begin{definition}\label{DefinitionNonDegenerate}
An irreducible subvariety $X$ of $\cA$ is said to be non-degenerate if
there exists an open non-empty subset $\Delta$ of $S^{\mathrm{an}}$,
with the Betti map
$b_{\Delta} \colon \cA_{\Delta}:=\pi^{-1}(\Delta) \rightarrow \mathbb{T}^{2g}$, 
such that
\begin{equation}\label{EqDefnNonDeg}
\max_{x \in \sman{X} \cap \cA_{\Delta}} \mathrm{rank}_{\IR} (\mathrm{d}b_{\Delta}|_{\sman{X}})_x = 2\dim X
\end{equation}
where $\mathrm{d}b_\Delta$ denotes the differential and $\sman{X}$ is
the analytification of the regular locus of $X$.
\end{definition}

As the inequality $\le$ in \eqref{EqDefnNonDeg} trivially holds true, \eqref{EqDefnNonDeg} is equivalent to: there exists $x \in \sman{X} \cap \cA_{\Delta}$ such that $\mathrm{rank}_{\IR} (\mathrm{d}b_{\Delta}|_{\sman{X}})_x = 2\dim X$.


In Proposition~\ref{PropBettiForm}(iii) we give another characterization of non-degenerate
subvarieties. We can now formulate the height inequality. 

\begin{theorem}\label{ThmHtInequality}
  Suppose that $\cA$ and $S$ are as above with $k=\IQbar$;
  in particular, $\cA$ satisfies {\tt (Hyp)}.
  Let $X$ be a closed irreducible subvariety of $\cA$
  defined over $\IQbar$ that dominates $S$.   Suppose $X$ is non-degenerate,
  as defined in Definition~\ref{DefinitionNonDegenerate}. Then there
  exist constants $c_1>0$ and $c_2\ge 0$ and a Zariski open dense subset
  $U$ of $X$ with 
  \begin{equation*}
    \hat{h}_{\cA}(P) \ge c_1 h( \pi(P) ) - c_2 \quad \text{for all}\quad P \in U(\IQbar).
  \end{equation*}
\end{theorem}


Note that \cite[Theorem~1.4]{GaoHab} is, up to some minor reduction,
precisely Theorem~\ref{ThmHtInequality} for $\dim S = 1$ together with
the criterion for $X$ to be non-degenerate when $\dim S = 1$. In general, the degeneracy behavior of $X$ is fully studied in \cite{GaoBettiRank}. See \cite[Theorem~1.1]{GaoBettiRank} for the criterion. However in practice, we sometimes still want to understand the height comparison on some degenerate $X$. One way to achieve this is by applying \cite[Theorem~1.3]{GaoBettiRank}, which asserts the following statement: If $X$ satisfies some reasonable properties, then we can apply Theorem~\ref{ThmHtInequality} after doing some simple operations with $X$. 

For the purpose of proving Proposition~\ref{PropAlgPtFar} and
furthermore Theorem~\ref{ThmBdFinRank}, we work in the following
situation. 



Let $\mathbb{A}_{g,\ell}$
denote the moduli space of principally polarized $g$-dimensional
abelian varieties with level-$\ell$-structure. It is
a classical fact that $\mathbb{A}_{g,\ell}$ is represented by an
 irreducible, regular, quasi-projective
variety defined over a number field,
see~\cite[Theorem~7.9 and below]{MFK:GIT94} or \cite[Theorem 1.9]{OortSteenbrink}, so it is a
fine moduli space.
% We will use the analytic description of $\mathbb{A}_{g,\ell}$ arising from a suitable  quotient of Siegel's upper half space.
Let $\mathbb{M}_{g,\ell}$ be the fine moduli space of smooth curves of genus
$g$ whose Jacobian is equipped with  level-$\ell$-structure;
see~\cite[(5.14)]{DM:irreducibility} or \cite[Theorem
1.8]{OortSteenbrink} for the existence.
Then $\mathbb{M}_{g,\ell}$ is an irreducible, regular, quasi-projective
variety defined over a number field. 
%for some integer $\ell\ge \bestlevel$, and
% let $\mathbb A_g$ be the fine moduli space of principally polarized abelian
% varieties of dimension $g$ with level-$\ell$-structure.

To avoid confusion on different notations in different references, we
make the following convention throughout  the paper. We will take
$\mathbb{A}_{g,\ell}$ and $\mathbb{M}_{g,\ell}$ as geometrically
irreducible varieties. 
%\Pcomment{I changed this to geometrically  irreducible as we sometimes treat them as defined over a number field.} 
 Some authors define
$\mathbb{A}_{g,\ell}$ over $\mathbb{Z}[1/\ell]$ (or over $\mathbb{Z}$)
and then consider it over $\IQbar$ by base change. The
$\IQbar$-variety thus obtained may not be irreducible, and each
irreducible component is defined over $\IQ(\zeta_{\ell})$ for some
 root of unity $\zeta_{\ell}$ of order $\ell$.
Choosing a geometrically irreducible component of $\mathbb{A}_{g,\ell}$
amounts to fixing a complex root of unity of order $\ell$.
We fix such a choice once and for all and consider
$\mathbb{A}_{g,\ell}$ as an irreducible variety defined over $\IQbar$.
% In our convention
% $\mathbb{A}_{g,\ell}$ will be an irreducible component of the
% $\IQbar$-variety thus obtained. Same
The same  holds
for $\mathbb{M}_{g,\ell}$.
We will usually fix $\ell$ and abbreviate $\mathbb{A}_{g,\ell}$ (resp.
$\mathbb{M}_{g,\ell}$) by $\mathbb{A}_g$ (resp. $\mathbb{M}_g$).
It is often convenient to consider $\mathbb{A}_g$ and $\mathbb{M}_g$
as over $\IQbar$, but sometimes we will recall that both arise from
varieties defined over the number field $\IQ(\zeta_\ell)$.
We denote the coarse moduli space of smooth curves of genus $g$
with $\mathbb{M}_{g,1}$.

Furthermore, let
$\mathfrak{C}_g \rightarrow \mathbb{M}_g$ be the universal curve and
$\mathfrak{A}_g \rightarrow \mathbb A_g$ be the universal abelian
variety. Taking the Jacobian of a smooth curve leads to the  Torelli morphism
$\mathbb{M}_g \rightarrow \mathbb A_g$ which is finite-to-$1$  (but
 not injective  as we
have level structure).
Moreover, for $M\ge 2$ let
$\mathscr{D}_M$ denote the $M$-th Faltings--Zhang morphism fiberwise defined by
sending
\begin{equation}
  \label{eq:faltingszhangintro}
  (P_0, P_1, \ldots, P_M) \mapsto (P_1-P_0, \ldots, P_M-P_0);
\end{equation}
we give a precise definition of this morphism
in $\mathsection$\ref{subsec:univfamily}. 
 % To ease notation we identify its image with $\mathbb{M}_g$.
Roughly speaking, we will  apply Theorem~\ref{ThmHtInequality} to
\[
  X := \mathscr{D}_M(\underbrace{\mathfrak{C}_g \times_{\mathbb{M}_g} \cdots \times_{\mathbb{M}_g} \mathfrak{C}_g}_{(M+1)\text{-copies}}) \subseteq \underbrace{\mathfrak{A}_g \times_{\mathbb M_g} \cdots \times_{\mathbb M_g} \mathfrak{A}_g}_{M\text{-copies}}
\]
for
% \Pcomment{Change right-hand side in display above
%   to fibered power of $\mathbb{M}_g$}
a suitable  $M$. To verify non-degeneracy we will refer to  
the second-named author's work \cite[Theorem~1.2']{GaoBettiRank} which
applies if $M$ is large in terms of $g$.  So we can apply Theorem~\ref{ThmHtInequality} to such $X$.  This
will eventually lead to Proposition~\ref{PropAlgPtFar}. %PropNTDistance}.

The morphism and its variants are powerful tools in diophantine
geometry, see~\cite[Lemma~4.1]{Faltings:DAAV}. It is closely connected
to problems involving small N\'eron--Tate height,
see~\cite[Lemma~3.1]{ZhangEquidist}% and\cite[Proposition~4.1]{DPvarabII}
. Stoll~\cite{Stoll:Uniform} used a
variant of \eqref{eq:faltingszhangintro} to show that a conjecture of
Pink~\cite{Pink} on unlikely intersections implies
Theorem~\ref{ThmBdFinRank} with the condition
$h(\iota([\mathrm{Jac}(C)])) \ge c_1$ removed and with $C$ allowed to
be defined over $\IC$.
 
At this stage it is worth outlining the main steps of the proof of \cite[Theorem~1.2']{GaoBettiRank}, or the more general \cite[Theorem~1.3]{GaoBettiRank}, due to its importance to the current paper. The major step is to establish a criterion, \textit{in simple geometric terms}, for an irreducible subvariety $X$ of the universal abelian variety $\mathfrak{A}_g$ to be degenerate. Roughly speaking, the proof of the desired criterion is divided into two steps. Step~1 transfers the degeneracy property to an \textit{unlikely intersection problem} in $\mathfrak{A}_g$ by invoking the \textit{mixed Ax--Schanuel theorem} for $\mathfrak{A}_g$ \cite[Theorem~1.1]{GaoMixedAS}. More precisely we show that $X$ is degenerate if and only if $X$ is the union of subvarieties satisfying an appropriate unlikely intersection property. Step~2 solves this unlikely intersection problem, and the key point is to use \cite[Theorem~1.4]{GaoMixedAS} to prove that the union mentioned above is a finite union. In this step the notion of \textit{weakly optimal subvarieties} introduced by the third-named author and Pila \cite{HabeggerPilaENS} is involved.


 %For $\pi^{\mathrm{univ}} \colon \mathfrak{A}_g \rightarrow \mathbb{A}_g$, the Betti map can be expressed explicitly using the Siegel upper half space $\mathfrak{H}_g$ (which is a universal covering of $\mathbb{A}_g^{\mathrm{an}}$) and an appropriate universal covering $\mathbf{u} \colon \mathcal X_{2g,\mathrm{a}} \rightarrow \mathfrak{A}_g^{\mathrm{an}}$. Let $\tilde{X}$ be a complex analytic irreducible component of $\mathbf{u}^{-1}(X^{\mathrm{sm}})$. If $X$ is degenerate, \textit{i.e.}, \eqref{EqDefnNonDeg} fails for $X$, then it is not hard to check that $\tilde{X}$ is covered by  holomorphic horizontal curves. Hence $X$ is union of $\mathbf{u}(\tilde{C})^{\mathrm{Zar}}$ where $\tilde{C}$ runs over all holomorphic horizontal curves contained in $\tilde{X}$. 
%After these preparations, two majors steps are required to obtain the desired criterion of degeneracy.  Step~1 is to study $\mathbf{u}(\tilde{C})^{\mathrm{Zar}}$. To do this, one applies the \textit{mixed Ax--Schanuel theorem} for $\mathfrak{A}_g$ \cite[Theorem~1.1]{GaoMixedAS} to $\tilde{C}$ and shows that the $\mathbf{u}(\tilde{C})^{\mathrm{Zar}}$'s thus obtained are precisely the subvarieties of $X$ which are characterized by a suitable \textit{unlikely intersection} property. Thus $X$ is degenerate if and only if $X$ satisfies a suitable unlikely intersection condition (more precisely, $X$ is the union of subvarieties satisfying a particular unlikely intersection property).  Step~2 is to study this unlikely intersection problem. The key point for this step is to prove that the union mentioned in Step~1 is a finite union in the following way. First one shows that the maximal elements in the union are \textit{weakly optimal} as defined by the third-named author and Pila \cite{HabeggerPilaENS}, then one invokes the second-named author's \cite[Theorem~1.4]{GaoMixedAS} to conclude that these maximal elements arise from finitely many families of subvarieties of $\mathfrak{A}_g$. With some extra standard computation on $\mathfrak{A}_g$ in terms of mixed Shimura varieties, one can then conclude. 
 
 
\subsection{General notation}
We collect here an overview of notation used throughout the text. 

Let $S$ be an irreducible, quasi-projective variety defined over an
algebraically closed field $k$. Then $\sm{S}$ denotes the regular
locus of $X$. If $\pi\colon \cA\rightarrow S$ is an abelian scheme
then $[N]\colon \cA\rightarrow\cA$ is the multiplication-by-$N$
morphism for all $N\in\IN=\{1,2,3,\ldots\}$, and if $s\in S(k)$, the fiber
$\cA_s=\pi^{-1}(s)$ is an abelian variety defined over $k$.
If $k\subset \IC$, then
$\an{S}$ denotes the analytification of $S$; it carries a natural
topology that is Hausdorff.

We write $\mathbb{T}$ for the circle group $\{z\in \IC : |z|=1\}$.

\subsection*{Acknowledgements} The authors would like to thank
Shou-wu Zhang for relevant discussions and Gabriel Dill for the
argument involving torsion points to bound $h_1$ on
page~\pageref{page:boundh1}. 
We would also like
to thank Lars K\"uhne and Ngaiming Mok  for discussions on Hermitian Geometry; our paper 
is much influenced by K\"uhne's approach towards bounded height
\cite{kuehne:semiabelianbhc}
and by Mok's approach to study the Mordell--Weil rank over function
fields \cite{Mok11Form}.
We thank Gabriel Dill, Lars K\"uhne, Fabien Pazuki, and Joseph H. Silverman for corrections and comments
on a draft of this paper. 
We also thank the referees for their careful reading and valuable comments. 
VD would like
to thank the NSF and the Giorgio and Elena Petronio Fellowship Fund II
for financial support for this work. VD has received funding from the European Union's Seventh Framework Programme (FP7/2007--2013) / ERC grant agreement n$^\circ$ 617129. ZG has received
funding from the French National Research Agency grant
ANR-19-ERC7-0004, and the European Research Council (ERC) under the
European Union's Horizon 2020 research and innovation programme (grant
agreement n$^\circ$ 945714). PH has received funding
from the Swiss National Science Foundation project n$^\circ$ 200020\_184623.
Both VD and ZG would like to
thank the Institute for Advanced Study and the special year ``Locally
Symmetric Spaces: Analytical and Topological Aspects'' for its
hospitality during this work.

%%% Local Variables:
%%% TeX-master: "main"
%%% End:

